\documentclass[a4paper,10pt]{report}\input{../0.Préambule/Préambule_cours_prof}\begin{document}

\newpage
\section*{Angles et parallélisme}
\addcontentsline{toc}{section}{Angles et parallélisme}

\begin{exer1}
Colorie d'une couleur différente chaque paire \\d'angles alternes-internes.

\vspace{-2.6cm}
\begin{flushright}
\begin{tikzpicture}[line cap=round,line join=round,>=triangle 45,x=1.6cm,y=1.0cm,scale = 0.5]
\clip(-1.,0.) rectangle (6.5,6.);
\draw [line width=1.2pt] (-0.48,1.4)-- (6.,2.);
\draw [line width=1.2pt] (-0.2,4.52)-- (6.,4.);
\draw [line width=1.2pt] (1.14,5.82)-- (5.12,-0.08);
\end{tikzpicture}
\end{flushright}
\end{exer1}

\vspace{-6mm}
\begin{exer1}
En t'aidant de la figure, complète les phrases

\begin{minipage}{6cm}
\includegraphics[scale=0.35]{../40.Image/12. Exo-Angles et parallélisme 1.png}
\end{minipage}
\begin{minipage}{10cm}
\begin{enumerate}[font=\color{exer1}]
\item $\widehat{rAt}$ et $\widehat{yBz}$ sont \dotfill \medskip
\item $\widehat{wAz}$ et $\widehat{zAr}$ sont \dotfill \medskip
\item \dotfill et $\widehat{wAB}$ sont alternes-internes. 
\end{enumerate}
\end{minipage}
\end{exer1}

\vspace{-2mm}
\begin{exer1}
On considère les angles déterminés par \\les droites $(EG)$ et $(AD)$


\vspace{-1cm}
\noindent \begin{minipage}{9cm}

\vspace{5mm}
Cite deux paires d'angles alternes-internes :
\begin{enumerate}[font=\color{exer1}]
\item déterminés par la sécante $(KC)$.
\item déterminés par la sécante $(BR)$.
\end{enumerate}
\end{minipage}
\begin{minipage}{6cm}
\includegraphics[scale=0.4]{../40.Image/12. Exo-Angles et parallélisme 2.png}
\end{minipage}
\end{exer1}

\vspace{-2mm}
\noindent \begin{minipage}{7cm}
\begin{exer1}
Colorie de la même couleur les angles de même mesure sachant que les droites $(AB)$ et $(CD)$ sont parallèles.

\includegraphics[scale=0.5]{../40.Image/12. Exo-Angles et parallélisme 3.png}
\end{exer1}
\end{minipage}
\hspace{0.8cm}
\begin{minipage}{9cm}
\begin{exer1}
Dans chaque cas, les droites $(d)$ et $(d')$ dont parallèles. Calcule mentalement puis écris la mesure de chaque angle, sans justifier.

\begin{center}
\includegraphics[scale=0.4]{../40.Image/12. Exo-Angles et parallélisme 4.png}
\end{center}
\end{exer1}
\end{minipage}

\vspace{-2mm}
\noindent \begin{minipage}{8cm}
\begin{exer2}
Les droites $(d')$ et $(d'')$ sont elles parallèles ? Justifie soigneusement.

\includegraphics[scale=0.4]{../40.Image/12. Exo-Angles et parallélisme 5.png}
\end{exer2}
\end{minipage}
\hspace{0.8cm}
\begin{minipage}{8cm}
\begin{exer2}
Les droites $(d')$ et $(d'')$ sont elles parallèles ? Justifie soigneusement.

\includegraphics[scale=0.4]{../40.Image/12. Exo-Angles et parallélisme 6.png}
\end{exer2}
\end{minipage}

\vspace{-2mm}
\noindent \begin{minipage}{6.5cm}
\begin{exer3}
Les points $A$, $D$ et $E$ sont alignés. Démontre que les droites $(AC)$ et $(DB)$ sont parallèles.
\begin{center}
\includegraphics[scale=0.33]{../40.Image/12. Exo-Angles et parallélisme 7.png}
\end{center}
\end{exer3}
\end{minipage}
\hspace{0.8cm}
\begin{minipage}{10cm}
\begin{exer3}
On considère la figure ci-contre.

\vspace{5mm}
\begin{minipage}{3.5cm}
\includegraphics[scale=0.4]{../40.Image/12. Exo-Angles et parallélisme 8.png}
\end{minipage}
\begin{minipage}{5.5cm}
\begin{enumerate}[font=\color{exer3}]
\item Démontre que $(NO)$ et $(LA)$ sont parallèles.
\item Démontre que les angles $\widehat{ALR}$ et $\widehat{NOR}$ ont la même mesure que tu calculeras.
\item Déduis-en la nature du triangle $NOR$.
\end{enumerate}
\end{minipage}
\end{exer3}
\end{minipage}

\begin{exer3}

\begin{enumerate}[font=\color{exer3}]
\item Construis une figure à main levée du parallélogramme $RIEN$ de centre $C$ tel que $CR=3$cm, $\widehat{CRI} = 35$\degre et $\widehat{CRN}$ est un angle droit. Tu indiqueras sur la figure la mesure des angles $\widehat{CEI}$ et $\widehat{CEN}$
\item Construis cette figure en vraie grandeur sans tracer 
\end{enumerate}
\end{exer3}

\end{document}
